% ============================================================================
% Chapter 43: The Future of Technology
% ============================================================================
\chapter{The Future of Technology}
\label{ch:future}

\begin{objectives}
  \item Explore emerging technologies shaping the future
  \item Understand the basics of AI, IoT, and cloud computing
  \item Think critically about technology's impact on society
  \item Get inspired to continue learning about computer science
\end{objectives}

\bytetip[fun]{You've reached the final chapter --- congratulations! Now let's look at what's coming next in the world of technology. Spoiler: it's AMAZING.}

\section{Emerging Technologies}
\label{sec:emerging}
\index{emerging technology}

% --- Figure 43.1: Future Tech Cards ---
\begin{figure}[H]
\centering
\begin{tabularx}{\textwidth}{@{}*{4}{>{\centering\arraybackslash}X}@{}}
  \begin{tcolorbox}[colback=FunSky!8, colframe=FunSky, width=3.4cm, height=3.5cm, valign=center, halign=center, arc=8pt]
    {\Huge\textcolor{FunSky}{\faRobot}}\vspace{6pt}\\
    \textbf{Artificial\\Intelligence}\vspace{3pt}\\
    {\scriptsize Machines that\\can learn and think}
  \end{tcolorbox}
  &
  \begin{tcolorbox}[colback=FunPurple!8, colframe=FunPurple, width=3.4cm, height=3.5cm, valign=center, halign=center, arc=8pt]
    {\Huge\textcolor{FunPurple}{\faVrCardboard}}\vspace{6pt}\\
    \textbf{Virtual\\Reality}\vspace{3pt}\\
    {\scriptsize Immersive digital\\worlds}
  \end{tcolorbox}
  &
  \begin{tcolorbox}[colback=LabGreen!8, colframe=LabGreen, width=3.4cm, height=3.5cm, valign=center, halign=center, arc=8pt]
    {\Huge\textcolor{LabGreen}{\faNetworkWired}}\vspace{6pt}\\
    \textbf{Internet of\\Things (IoT)}\vspace{3pt}\\
    {\scriptsize Smart devices\\connected online}
  \end{tcolorbox}
  &
  \begin{tcolorbox}[colback=FunCoral!8, colframe=FunCoral, width=3.4cm, height=3.5cm, valign=center, halign=center, arc=8pt]
    {\Huge\textcolor{FunCoral}{\faCloud}}\vspace{6pt}\\
    \textbf{Cloud\\Computing}\vspace{3pt}\\
    {\scriptsize Services delivered\\over the internet}
  \end{tcolorbox}
\end{tabularx}
\caption{Four technologies shaping our future}
\label{fig:future-tech}
\end{figure}

\section{Artificial Intelligence (AI)}
\label{sec:ai}
\index{artificial intelligence}

\hl{Artificial Intelligence} is technology that enables machines to learn from data, recognise patterns, and make decisions --- almost like a human brain, but digital.

\begin{realworld}
You already use AI every day! When YouTube recommends videos, when your phone recognises your face, when Google completes your search query, and when Siri or Alexa answers your questions --- that's all AI at work!
\end{realworld}

\section{Internet of Things (IoT)}
\label{sec:iot}
\index{Internet of Things}

The \hl{Internet of Things} connects everyday objects to the internet --- from smart watches and fitness trackers to smart fridges and home assistants.

\begin{table}[H]
\centering
\caption{Examples of IoT devices}
\label{tab:iot}
\renewcommand{\arraystretch}{1.3}
\begin{tabularx}{\textwidth}{l X X}
\toprule
\rowcolor{HeaderRow}
\textcolor{white}{\textbf{Device}} & \textcolor{white}{\textbf{What It Does}} & \textcolor{white}{\textbf{How It Helps}} \\
\midrule
\rowcolor{RowLight} Smart Watch & Tracks health, shows notifications & Monitors heart rate, steps \\
\rowcolor{RowDark} Smart Speaker & Plays music, answers questions & Voice-controlled assistance \\
\rowcolor{RowLight} Smart Thermostat & Controls home temperature & Saves energy automatically \\
\rowcolor{RowDark} Smart Lock & Locks/unlocks doors remotely & Security from your phone \\
\bottomrule
\end{tabularx}
\end{table}

\section{Your Journey Continues!}
\label{sec:journey}

This book has given you a strong foundation in computer science. But this is just the \textbf{beginning}! Here are some exciting paths you can explore next:

\begin{itemize}[leftmargin=*, label={\textcolor{ModuleBlue}{\faRocket}}, itemsep=4pt]
  \item \textbf{Programming} --- Learn Python, Scratch, or JavaScript to create your own programs
  \item \textbf{Web Development} --- Build your own websites with HTML, CSS, and JavaScript
  \item \textbf{Robotics} --- Build and program real robots
  \item \textbf{Cybersecurity} --- Protect systems from hackers and threats
  \item \textbf{Game Development} --- Create your own video games
  \item \textbf{Data Science} --- Analyse data to discover patterns and insights
\end{itemize}

\bytetip[tip]{The most important skill in computer science isn't memorising facts --- it's \textbf{curiosity}. Keep asking ``How does this work?'' and ``What if I tried this?'' That's how every great inventor started!}

\begin{funfact}
The term ``Artificial Intelligence'' was coined in 1956 by \textbf{John McCarthy} at a conference at Dartmouth College. He predicted that machines would be as smart as humans within 20 years. It took a bit longer than that, but we're getting closer every day!
\end{funfact}

\begin{reviewqs}
  \item What is Artificial Intelligence? Give two examples of AI you use daily.
  \item What does IoT stand for? Name three IoT devices.
  \item What is cloud computing, and name two cloud services.
  \item How might AI change education in the next 10 years?
  \item What area of computer science interests you most, and why?
\end{reviewqs}

\begin{challenge}
Write a \textbf{one-page essay} titled ``The Computer of 2050'' describing what you think computers will look like and be able to do in 25 years. Will they be wearable? Will they read your mind? Use your imagination --- and back it up with real technology trends!
\end{challenge}

\vspace{20pt}
\begin{center}
\begin{tcolorbox}[
  enhanced, colback=FunYellow!15, colframe=FunYellow!80!black, arc=12pt,
  width=10cm, halign=center,
  fonttitle=\Large\bfseries\sffamily, title={\faStar\ Congratulations!},
  coltitle=white, boxed title style={colback=FunYellow!80!black}
]
\large\sffamily
You've completed \textbf{Basic computer course !} \\[8pt]
You now have the skills to use computers confidently, create professional documents, and stay safe online.\\[8pt]
\textbf{Keep exploring. Keep creating. Keep learning!}\\[6pt]
\normalsize\itshape --- Byte the Robot \faRobot
\end{tcolorbox}
\end{center}
